% !TEX program = xelatex
% 第三讲：laplace_transform.tex + inverse_laplace_transform.tex + applications.tex
\documentclass[fontset=windows, chinese, 10pt, xcolor=dvipsnames]{beamer}
% Shared preamble for the chapter 04 standalone lectures.
% Each lecture sets \title and \subtitle after this file is read.

\usepackage{amsmath,amssymb,mathtools}
\usepackage{xeCJK}
\usepackage{zhnumber}
\usepackage{tikz}
\usetikzlibrary{arrows.meta,calc,positioning}

\usetheme{metropolis}
\usepackage{beamercolorthemeHistNanjing}

\newcommand{\red}[1]{\textcolor{red}{#1}}
\newcommand{\blue}[1]{\textcolor{blue}{#1}}
\newcommand{\green}[1]{\textcolor{green}{#1}}
\newcommand{\gray}[1]{\textcolor{gray}{#1}}
\newcommand{\orange}[1]{\textcolor{orange}{#1}}

\newcommand{\sssub}[1]{ {\mathrm{#1} }}  % script size subscript\\\
\newcommand{\funcd}[2]{ \frac{d{#1}}{d{#2}}} %function full derivative
\newcommand{\funcpd}[2]{ \frac{\partial{#1}}{\partial{#2}}} %function partial derivative
\newcommand{\fnald}[2] { \frac{\delta {#1}}{\delta {#2}}} % functional derivative

\renewcommand{\Re}{\operatorname{Re}}
\renewcommand{\Im}{\operatorname{Im}}
\def\ext{\sssub{ext}}
\def\x{\sssub{x}}
\def\xc{\sssub{xc}}
\def\H{\sssub{H}}
\def\Hx{\sssub{Hx}}
\def\Hxc{\sssub{Hxc}}
\def\c{\sssub{c}}
\def\s{\sssub{s}}

% \def\ext{\sssub{ext}}
% \def\ext{\sssub{ext}}
% \def\ext{\sssub{ext}}
% \def\ext{\sssub{ext}}
% \defsssub{ext}


% definitions of bold faced symbols
\def\brpppp{{\mathbf{r}^{\prime\prime\prime\prime}}}
\def\brppp{{\mathbf{r}^{\prime\prime\prime}}}
\def\brpp{{\mathbf{r}^{\prime\prime}}}
\def\brp{{\mathbf{r}^{\prime}}}
\def\bzp{{\mathbf{z}^{\prime}}}
\def\bxp{{\mathbf{x}^{\prime}}}
\def\tp{{{t}^{\prime}}}
\def\tpp{{{t}^{\prime\prime}}}
\def\tppp{{{t}^{\prime\prime\prime}}}

\def\tbr{{\tilde{\mathbf{r}}}}
\def\bk{{\mathbf{k}}}
\def\br{{\mathbf{r}}}
\def\bp{{\mathbf{p}}}
\def\bq{{\mathbf{q}}}
\def\bv{{\mathbf{v}}}
\def\bz{{\mathbf{z}}}
\def\mbf{{\mathbf{f}}}

\def\bx{{\mathbf{x}}}
\def\bR{{\mathbf{R}}}
\def\bM{{\mathbf{M}}}
\def\bP{{\mathbf{P}}}
\def\bT{{\mathbf{T}}}
\def\bK{{\mathbf{K}}}
\def\bA{{\mathbf{A}}}
\def\bB{{\mathbf{B}}}
\def\bD{{\mathbf{D}}}
\def\bE{{\mathbf{E}}}
\def\bF{{\mathbf{F}}}
\def\bG{{\mathbf{G}}}
\def\bH{{\mathbf{H}}}
\def\bI{{\mathbf{I}}}
\def\bJ{{\mathbf{J}}}
\def\bP{{\mathbf{P}}}
\def\bS{{\mathbf{S}}}
\def\bV{{\mathbf{V}}}
\def\bX{{\mathbf{X}}}
\def\bY{{\mathbf{Y}}}
\def\dd{{\mathop{}\!\mathrm{d}}}
\def\dr{{\!\!\dd\br\,}}
\def\dx{{\!\!\dd\bx\,}}
\def\dy{{\!\!\dd\by\,}}
\def\dk{{\!\!\dd\bk\,}}
\def\dG{{\!\!\dd\bG\,}}


\def\half{{\frac{1}{2}}}
\def\tr{{\mathrm{tr}}}


\def\Arg{{\mathrm{Arg}\;}}
\def\arg{{\mathrm{arg}\;}}
\def\Res{{\mathrm{Res}\;}}


\author{罗凯}
\institute{南京理工大学物理学院}
\date{\zhdigits{\the\year},\zhganzhinian{\the\year}年秋}

\usebackgroundtemplate{%
\begin{tikzpicture}[remember picture,overlay]
  \node[opacity=0.05, at=(current page.center)] {
    \includegraphics[width=0.6\paperwidth]{njust-logo.jpeg}
  };
\end{tikzpicture}}


\title{傅里叶分析}
\subtitle{第三讲：拉普拉斯变换}

\begin{document}

\begin{frame}
  \titlepage
\end{frame}

\begin{frame}{傅里叶变换还缺什么}
  傅里叶变换适合看频谱，用在初值问题上有两处限制。
  \begin{itemize}
    \item 绝对可积是充分条件。$e^{\alpha t}$（$\alpha>0$）在 $t\to+\infty$ 时发散，不能直接套定义。
    \item 反变换是对实频率的积分，常常不好算。
  \end{itemize}
  拉普拉斯变换只关心 $t\geq 0$，并乘上衰减因子 $e^{-\sigma t}$，把问题收成右半平面上的复变函数。
\end{frame}

\begin{frame}{本讲提纲}
  \begin{enumerate}
    \item 由因果函数和收敛因子写出拉普拉斯变换，并回到 Bromwich 积分；
    \item 基本像函数：常数、幂、指数、双曲函数和三角函数；
    \item 导数定理与卷积，以及用留数、裂项、延迟求原函数；
    \item 三个应用：级数、含 $1/t$ 的积分、简谐振子。
  \end{enumerate}
\end{frame}

% =============================================================
\section{定义}

\begin{frame}{只保留 $t\geq 0$}
  初值问题已知 $f(0)$，求 $t>0$ 的变化。$t<0$ 的值置为零：
  \begin{equation*}
    f(t)=0\qquad(t<0).
  \end{equation*}
  选 $\sigma>0$，令 $g(t)=e^{-\sigma t}f(t)$。若 $|f(t)|$ 不超过某个指数，$g$ 就绝对可积，可以对它做傅里叶变换。
\end{frame}

\begin{frame}{定义}
  对 $g$ 作变换，并记 $p=\sigma+\mathrm{i}\omega$。定义
  \begin{equation*}
    \boxed{\bar f(p)=\mathcal{L}[f]
      =\int_0^{\infty} e^{-pt}f(t)\,\mathrm{d}t.}
  \end{equation*}
  这个积分是拉普拉斯积分，$e^{-pt}$ 是核。积分从 $0$ 开始，不是从 $-\infty$。
  \vspace{0.4em}

  $\bar f$ 是像函数，$f$ 是原函数。变换是线性的：
  \begin{equation*}
    \mathcal{L}[af+bg]=a\mathcal{L}[f]+b\mathcal{L}[g].
  \end{equation*}
  傅里叶变换同样是线性的。
\end{frame}

\begin{frame}{反变换：Bromwich 积分}
  把 $g$ 的傅里叶逆变换写回 $f$，再用 $p=\sigma+\mathrm{i}\omega$，$\mathrm{d}\omega=\mathrm{d}p/\mathrm{i}$，得到
  \begin{equation*}
    f(t)=\frac{1}{2\pi\mathrm{i}}
      \int_{\sigma-\mathrm{i}\infty}^{\sigma+\mathrm{i}\infty}
      \bar f(p)\,e^{pt}\,\mathrm{d}p.
  \end{equation*}
  积分路径是竖直线 $\operatorname{Re}p=\sigma$，位于像函数的解析区域里。这就是 Bromwich 积分，也叫 Riemann--Mellin 积分。
\end{frame}

\begin{frame}{存在的充分条件}
  \begin{enumerate}
    \item 在每个有限区间 $[0,T]$ 上，$f$ 只有有限个间断点，其余点 $f$ 及其导数连续；
    \item 存在 $M>0$ 和 $\sigma_0\geq 0$，使 $|f(t)|<M e^{\sigma_0 t}$。
  \end{enumerate}
  $\sigma_0$ 的下界叫收敛横标。$\operatorname{Re}p>\sigma_0$ 时积分收敛，并且 $\bar f(p)$ 在这个半平面里解析。
  \vspace{0.3em}

  $e^{t^2}$ 比任何指数都长得快，不满足第二条。满足条件时，$\operatorname{Re}p\to+\infty$ 必有 $\bar f(p)\to 0$。像函数若在无穷远处像 $p$ 的正幂次那样变大，它就不是一个拉普拉斯像。
\end{frame}

\begin{frame}{$\delta$ 函数的像}
  $t_0>0$ 时，挑选性给出
  \begin{equation*}
    \mathcal{L}[\delta(t-t_0)]
    =\int_0^{\infty} e^{-pt}\delta(t-t_0)\,\mathrm{d}t
    =e^{-pt_0}.
  \end{equation*}
  延迟一个脉冲，像函数就乘上 $e^{-pt_0}$。这和后面的延迟定理是同一件事。
\end{frame}

% =============================================================
\section{基本像函数}

\begin{frame}{例：常数与 $t$}
  $\operatorname{Re}p>0$ 时
  \begin{equation*}
    \mathcal{L}[1]=\int_0^{\infty} e^{-pt}\,\mathrm{d}t=\frac{1}{p}.
  \end{equation*}
  分部积分，$\displaystyle\int_0^{\infty} t e^{-pt}\,\mathrm{d}t=1/p^2$。
  继续做下去，
  \begin{equation*}
    \mathcal{L}[t^n]=\frac{n!}{p^{n+1}},
    \qquad n=0,1,2,\ldots.
  \end{equation*}
  检查：$p\to\infty$ 时这些像都趋于 $0$。
\end{frame}

\begin{frame}{例：指数}
  $\operatorname{Re}(p-s)>0$ 时
  \begin{equation*}
    \mathcal{L}[e^{st}]=\int_0^{\infty} e^{-(p-s)t}\,\mathrm{d}t=\frac{1}{p-s}.
  \end{equation*}
  再乘上 $t$，仍用分部积分，
  \begin{equation*}
    \mathcal{L}[t^n e^{st}]=\frac{n!}{(p-s)^{n+1}}.
  \end{equation*}
  $s$ 把收敛横标从 $0$ 挪到 $\operatorname{Re}s$。$s$ 为正时，原函数指数增长，像函数的极点在 $p=s$。
\end{frame}

\begin{frame}{例：双曲函数}
  $\cosh(kt)=\tfrac12(e^{kt}+e^{-kt})$，$\sinh(kt)=\tfrac12(e^{kt}-e^{-kt})$。用上一页的线性组合，
  \begin{align*}
    \mathcal{L}[\cosh(kt)]
      &=\frac{p}{p^2-k^2},\\
    \mathcal{L}[\sinh(kt)]
      &=\frac{k}{p^2-k^2}.
  \end{align*}
  要求 $\operatorname{Re}p>|k|$，两条指数才都收敛。
\end{frame}

\begin{frame}{正弦和余弦}
  $\displaystyle\int_0^{\infty} e^{\mathrm{i}\omega t}e^{-pt}\,\mathrm{d}t
    =1/(p-\mathrm{i}\omega)$，在 $\operatorname{Re}p>0$ 时成立。
  取虚部和实部：
  \begin{align*}
    \mathcal{L}[\sin(\omega t)]&=\frac{\omega}{p^2+\omega^2},\\
    \mathcal{L}[\cos(\omega t)]&=\frac{p}{p^2+\omega^2}.
  \end{align*}
  把双曲函数里的 $k$ 换成 $\mathrm{i}\omega$，分母 $p^2-k^2$ 变成 $p^2+\omega^2$，得到同一对公式。简谐振子的像就是余弦这一条。
\end{frame}

% =============================================================
\section{性质}

\begin{frame}{导数定理}
  分部积分，并假定 $e^{-pt}f(t)$ 在无穷远处消失：
  \begin{equation*}
    \int_0^{\infty} f'(t)\,e^{-pt}\,\mathrm{d}t
    =-f(0)+p\int_0^{\infty} f(t)\,e^{-pt}\,\mathrm{d}t.
  \end{equation*}
  因此
  \begin{equation*}
    \boxed{\mathcal{L}[f']=p\bar f(p)-f(0).}
  \end{equation*}
  和傅里叶变换的 $\mathrm{i}\omega F$ 相比，多出来的 $-f(0)$ 正是初值。这是拉普拉斯变换能解初值问题的原因。
\end{frame}

\begin{frame}{高阶导数}
  对 $f'$ 再使用一次导数定理，
  \begin{equation*}
    \mathcal{L}[f'']=p^2\bar f(p)-pf(0)-f'(0).
  \end{equation*}
  一般地，
  \begin{equation*}
    \begin{aligned}
    \mathcal{L}[f^{(n)}]
      =p^n\bar f(p)
      &-p^{n-1}f(0)-p^{n-2}f'(0)\\
      &-\cdots-f^{(n-1)}(0).
    \end{aligned}
  \end{equation*}
  二阶方程会用到 $f(0)$ 和 $f'(0)$，不多也不少。
\end{frame}

\begin{frame}{积分、相似、位移、延迟}
  \begin{align*}
    \mathcal{L}\Bigl[\int_0^t \psi(\tau)\,\mathrm{d}\tau\Bigr]
      &=\frac{1}{p}\mathcal{L}[\psi],\\
    \mathcal{L}[f(at)]
      &=\frac{1}{a}\bar f\Bigl(\frac{p}{a}\Bigr)
      \quad(a>0),\\
    \mathcal{L}[e^{-\lambda t}f(t)]
      &=\bar f(p+\lambda),\\
    \mathcal{L}[f(t-t_0)]
      &=e^{-pt_0}\bar f(p).
  \end{align*}
  延迟时要记住 $t<t_0$ 的原函数是 $0$：平移以前，函数在负半轴上已经被截断。
\end{frame}

\begin{frame}{卷积}
  \begin{equation*}
    (f_1*f_2)(t)
    =\int_0^t f_1(\tau)\,f_2(t-\tau)\,\mathrm{d}\tau.
  \end{equation*}
  上限是 $t$，因为两个原函数在负时间都是零。此时没有额外的 $2\pi$：
  \begin{equation*}
    \boxed{\mathcal{L}[f_1*f_2]=\bar f_1(p)\,\bar f_2(p).}
  \end{equation*}
  像函数相乘，对应原函数从 $0$ 卷积到 $t$。
\end{frame}

% =============================================================
\section{反变换}

\begin{frame}{留数}
  $t>0$ 时，$e^{pt}$ 在左半平面衰减。若 $\bar f(p)$ 在左方的大圆弧上一致趋于零，推广的若尔当引理说，那段弧上的积分趋于零。
  \vspace{0.4em}

  竖直线左边的奇点因此都被回路圈进来。原函数等于这些奇点上的留数之和：
  \begin{equation*}
    f(t)=\sum_j
      \operatorname{Res}\bigl[\bar f(p)\,e^{pt},\,p_j\bigr].
  \end{equation*}
  右半平面在收敛横标的右边没有奇点，不用加进去。
\end{frame}

\begin{frame}{像函数对 $p$ 求导}
  在 $\operatorname{Re}p>\sigma_0$ 内可以在积分号下求导：
  \begin{equation*}
    \frac{\mathrm{d}^n}{\mathrm{d}p^n}\bar f(p)
    =\int_0^{\infty}(-t)^n f(t)\,e^{-pt}\,\mathrm{d}t.
  \end{equation*}
  所以 $\bar f^{(n)}(p)$ 的原函数是 $(-t)^n f(t)$。
  \vspace{0.3em}

  例如 $\mathcal{L}[1]=1/p$，对 $p$ 求导一次，
  \begin{equation*}
    \frac{1}{p^2}=-\frac{\mathrm{d}}{\mathrm{d}p}\frac{1}{p}
  \end{equation*}
  的原函数是 $t$。这和直接计算 $\mathcal{L}[t]$ 一致。
\end{frame}

\begin{frame}{例：有理函数裂项}
  \begin{equation*}
    \frac{1}{p^3(p+\alpha)}
    =\frac{1}{\alpha}\frac{1}{p^3}
     -\frac{1}{\alpha^2}\frac{1}{p^2}
     +\frac{1}{\alpha^3}\frac{1}{p}
     -\frac{1}{\alpha^3}\frac{1}{p+\alpha}.
  \end{equation*}
  令 $p=0$ 得到 $1/p^3$ 的系数，令 $p=-\alpha$ 得到最后一项。逐项反演：
  \begin{equation*}
    \frac{t^2}{2\alpha}-\frac{t}{\alpha^2}
      +\frac{1}{\alpha^3}
      -\frac{e^{-\alpha t}}{\alpha^3}.
  \end{equation*}
  $1/p^2$ 的系数是负的，所以 $t$ 的系数是 $-1/\alpha^2$。
\end{frame}

\begin{frame}{支点：$1/\sqrt{p}$}
  $p=0$ 是支点，不能只画一条竖直线再补左半圆弧。从 $0$ 沿负实轴切开，辐角在割缝上沿取 $\pi$，下沿取 $-\pi$。
  \vspace{0.3em}

  大圆弧和小圆的贡献都趋于零，闭路内没有奇点。竖直线上的积分于是等于割缝两侧的贡献。令 $y=\sqrt{|p|t}$，
  \begin{equation*}
    \mathcal{L}^{-1}\bigl[p^{-1/2}\bigr]
    =\frac{1}{\sqrt{\pi t}}.
  \end{equation*}
  反过来说，$\mathcal{L}[1/\sqrt{\pi t}]=1/\sqrt{p}$。
\end{frame}

\begin{frame}{割缝两侧怎么加}
  上沿 $\sqrt{p}=\mathrm{i}\sqrt{r}$，下沿 $\sqrt{p}=-\mathrm{i}\sqrt{r}$。两侧相加、再除以 $2\pi\mathrm{i}$ 之后，
  \begin{equation*}
    f(t)=\frac{1}{\pi}\int_0^{\infty}\frac{e^{-rt}}{\sqrt{r}}\,\mathrm{d}r.
  \end{equation*}
  令 $y=\sqrt{rt}$，则 $\mathrm{d}r=2y\,\mathrm{d}y/t$。积分化为
  \begin{equation*}
    f(t)=\frac{2}{\pi\sqrt{t}}\int_0^{\infty}e^{-y^2}\,\mathrm{d}y
    =\frac{1}{\sqrt{\pi t}},
  \end{equation*}
  其中 $\displaystyle\int_0^{\infty}e^{-y^2}\,\mathrm{d}y=\sqrt{\pi}/2$。
\end{frame}

\begin{frame}{例：延迟与卷积一起用}
  求 $e^{-\alpha p}/[p(p+b)]$ 的原函数，$b\neq 0$。
  \vspace{0.3em}

  $\mathcal{L}[\Theta(t)]=1/p$，延迟定理给出 $e^{-\alpha p}/p$ 的原函数 $\Theta(t-\alpha)$。
  $1/(p+b)$ 的原函数是 $e^{-bt}$。二者相乘对应卷积：
  \begin{align*}
    \int_0^t \Theta(\tau-\alpha)\,e^{-b(t-\tau)}\,\mathrm{d}\tau
    &=\Theta(t-\alpha)\int_{\alpha}^{t} e^{-b(t-\tau)}\,\mathrm{d}\tau\\
    &=\frac{1-e^{-b(t-\alpha)}}{b}\,\Theta(t-\alpha).
  \end{align*}
  $t<\alpha$ 时结果是零：信号还没有到达。
\end{frame}

% =============================================================
\section{应用}

\begin{frame}{级数 $\sum 1/n^2$}
  由 $\displaystyle\int_0^{\infty} t e^{-pt}\,\mathrm{d}t=1/p^2$，
  \begin{align*}
    \sum_{n=1}^{\infty}\frac{1}{n^2}
    &=\sum_{n=1}^{\infty}\int_0^{\infty} t\,e^{-nt}\,\mathrm{d}t\\
    &=\int_0^{\infty}\frac{t}{e^{t}-1}\,\mathrm{d}t.
  \end{align*}
  交换求和与积分是因为被积函数非负。这个积分等于 $\pi^2/6$，和第一讲用 $f(x)=x$ 的帕塞瓦尔等式得到的结果相同。
\end{frame}

\begin{frame}{含 $1/t$ 的积分}
  若 $\displaystyle\int_p^{\infty}\bar f(q)\,\mathrm{d}q$ 存在，并且 $t\to 0$ 时 $|f(t)/t|$ 有界，则
  \begin{equation*}
    \frac{f(t)}{t}
    \quad\text{的像函数是}\quad
    \int_p^{\infty}\bar f(q)\,\mathrm{d}q.
  \end{equation*}
  再令 $p\to 0$，
  \begin{equation*}
    \int_0^{\infty}\frac{f(t)}{t}\,\mathrm{d}t
    =\int_0^{\infty}\bar f(p)\,\mathrm{d}p.
  \end{equation*}
\end{frame}

\begin{frame}{例：$\displaystyle\int_0^{\infty}(\sin t)/t\,\mathrm{d}t$}
  $\mathcal{L}[\sin t]=1/(p^2+1)$。于是
  \begin{equation*}
    \int_0^{\infty}\frac{\sin t}{t}\,\mathrm{d}t
    =\int_0^{\infty}\frac{1}{p^2+1}\,\mathrm{d}p
    =\frac{\pi}{2}.
  \end{equation*}
  被积函数在 $t=0$ 的极限是 $1$，奇点是可去的。直接在实轴上用留数并不方便，像函数的积分是反正切。
\end{frame}

\begin{frame}{例：$\displaystyle\int_0^{\infty}\dfrac{\cos(at)-\cos(bt)}{t}\,\mathrm{d}t$}
  $a>0$，$b>0$。
  $\mathcal{L}[\cos(at)-\cos(bt)]=p/(p^2+a^2)-p/(p^2+b^2)$，于是
  \begin{align*}
    \int_0^{\infty}\frac{\cos(at)-\cos(bt)}{t}\,\mathrm{d}t
    &=\int_0^{\infty}
      \Bigl(\frac{p}{p^2+a^2}-\frac{p}{p^2+b^2}\Bigr)\mathrm{d}p\\
    &=\frac12\Bigl[\ln\frac{p^2+a^2}{p^2+b^2}\Bigr]_0^{\infty}
    =\ln\frac{b}{a}.
  \end{align*}
  上限趋于 $\ln 1=0$，下限留下 $-\ln(a/b)$。$t=0$ 处的极限是 $0$，同样是可去奇点。
\end{frame}

\begin{frame}{例：简谐振子}
  \begin{equation*}
    m X''+k X=0,
    \qquad X(0)=X_0,\quad X'(0)=0.
  \end{equation*}
  设 $x(p)=\mathcal{L}[X]$。导数定理把方程变成代数方程
  \begin{equation*}
    m\bigl(p^2 x(p)-p X_0\bigr)+k x(p)=0.
  \end{equation*}
  令 $\omega_0^2=k/m$，
  \begin{equation*}
    x(p)=X_0\frac{p}{p^2+\omega_0^2}.
  \end{equation*}
  这是 $\cos(\omega_0 t)$ 的像，所以 $X(t)=X_0\cos(\omega_0 t)$。
  初速度为零，解里没有正弦项。
\end{frame}

\begin{frame}{若初速度不是零}
  设 $X'(0)=V_0$。二阶导数定理多一项 $-m V_0$：
  \begin{equation*}
    x(p)=X_0\frac{p}{p^2+\omega_0^2}
        +V_0\frac{1}{p^2+\omega_0^2}.
  \end{equation*}
  第二项是 $\sin(\omega_0 t)/\omega_0$ 的像，因此
  \begin{equation*}
    X(t)=X_0\cos(\omega_0 t)+\frac{V_0}{\omega_0}\sin(\omega_0 t).
  \end{equation*}
  微分方程变成了裂项。这就是本讲把导数定理放在卷积和反演前面的原因。
\end{frame}

\begin{frame}{本讲小结}
  \begin{itemize}
    \item $\mathcal{L}[f]=\int_0^{\infty} e^{-pt}f(t)\,\mathrm{d}t$。它允许指数增长，并把 $t<0$ 截掉。
    \item $\mathcal{L}[f']=p\bar f-f(0)$。初值出现在像函数里，所以振子方程可以直接解。
    \item 有理像函数用裂项和留数；支点要切开。卷积与延迟用来处理 $e^{-\alpha p}$ 这种因子。
  \end{itemize}
\end{frame}

\end{document}
